%% ==================================================================
%% Draft sections for: The Kähler Substrate: A Multi-Layer Geometric
%% Engine for the Davis Framework, Validated by Sheaf-Composition
%% Runtime Observation
%%
%% Status: DRAFT v0.1, sections built up one at a time for bee's
%% review. Preamble below matches the sheaf-composition paper's
%% template; sections will be added between \begin{document} and
%% \end{document} in the drafting-plan order.
%%
%% Drafting order (per outline §"Drafting plan"):
%%   §5 → §6 → §3 → §4 → §7 → §8 → §1 + §2 + §9 → abstract
%%
%% Drafting standard (bee's instruction, 2026-05-25):
%%   1. First-reader onboarding via inline \paragraph{Intuition.}
%%   2. Numbered lemmas + proofs behind every load-bearing claim.
%%   3. Real-data tables backing every numeric claim.
%% ==================================================================

\documentclass[11pt,letterpaper]{article}
\usepackage[utf8]{inputenc}
\usepackage[T1]{fontenc}
\usepackage{lmodern}
\usepackage{textcomp}
\usepackage[margin=1in]{geometry}
\emergencystretch=3em
\usepackage{amsmath,amssymb,amsthm}
\usepackage{mathtools}
\usepackage{hyperref}
\usepackage{booktabs}
\usepackage{multirow}
\usepackage{float}
\usepackage{graphicx}
\usepackage{enumitem}
\usepackage{listings}
\usepackage{xcolor}

\newtheorem{theorem}{Theorem}[section]
\newtheorem{lemma}[theorem]{Lemma}
\newtheorem{proposition}[theorem]{Proposition}
\newtheorem{corollary}[theorem]{Corollary}
\newtheorem{definition}[theorem]{Definition}
\newtheorem{remark}[theorem]{Remark}

\hypersetup{
  colorlinks=true,
  linkcolor=black,
  urlcolor=blue,
  citecolor=black,
}

\lstdefinestyle{gql}{
  basicstyle=\ttfamily\small,
  breaklines=true,
  frame=single,
  framesep=4pt,
  keywordstyle=\bfseries,
}
\lstdefinestyle{algo}{
  basicstyle=\ttfamily\small,
  breaklines=true,
  frame=single,
  framesep=4pt,
}

\title{\textbf{The K\"ahler Substrate} \\[6pt]
\large A Multi-Layer Geometric Engine for the Davis Framework, \\
Validated by Sheaf-Composition Runtime Observation \\[10pt]
\normalsize \emph{Draft v0.1 --- partial: §5 only}}

\author{A. Bee Rosa Davis \\
\href{mailto:bee_davis@alumni.brown.edu}{\nolinkurl{bee_davis@alumni.brown.edu}} \\
Independent Researcher}

\date{\today}

\begin{document}
\maketitle

\noindent\fbox{\parbox{0.97\linewidth}{\small
\textbf{Reading note for this draft.} This PDF contains only §5
(``The optionality contract --- a load-bearing engineering
property''). It is shipped as a section-shape rehearsal so bee can
read the prose register, table style, lemma form, and first-reader
onboarding rhythm before the rest of the paper is drafted. The
remaining sections (§§1--4, 6--9 plus appendices) will land
incrementally per the drafting plan in
\texttt{theory/kahler\_upgrade/PAPER\_OUTLINE\_KAHLER\_SUBSTRATE\_v0.1.md}.

\medskip

\textbf{Two display quirks of the rehearsal-only PDF.}
(1) Section and subsection numbers display as ``1, 1.1, 1.2, \dots''
because this is the only section in the document; when §1--§9 are
assembled the numbers will read ``5, 5.1, 5.2, \dots'' as the
outline specifies.
(2) Cross-references to other sections (e.g. ``\S\ref{sec:layers}''
in the captions of Tables~\ref{tab:python-validation}) appear as
``\textbf{??}'' until those sections exist; that is expected.}}

\bigskip



%% Paper body sections will land here per the v0.2 drafting plan.
%% Currently empty awaiting §4 (the prediction) draft.

\end{document}
